\SourceLeaf{074}
\[
\begin{aligned}&\frac{(ax+by+cz+\text{etc.})^2}{p}+\frac{(a'x+b'y+c'z+\text{etc.})^2}{p'}\\&\qquad+\frac{(a''x+b''y+c''z+\text{etc.})^2}{p''}+\text{etc.}\quad(15)\end{aligned}
\]

Porro patet. esse
\[
T=x\xi+y\eta+z\zeta+\text{etc.}\qquad(16)
\]
et si hic denuo \(x,y,z\) etc. adiumento aequationum (7) per \(\xi,\eta,\zeta\) etc. exprimuntur,
\[
\begin{aligned}
T={}&[\alpha\alpha]\xi\xi+2[\alpha\beta]\xi\eta+2[\alpha\gamma]\xi\zeta+\text{etc.}\\&+[\beta\beta]\eta\eta+2[\beta\gamma]\eta\zeta+\text{etc.}\\
&+[\gamma\gamma]\zeta\zeta+\text{etc.}
\end{aligned}
\]

Theoria supra evoluta bina systemata valorum determinatorum quantitatum \(x,y,z\) etc., atque \(\xi,\eta,\zeta\) etc. continet; priori, in quo \(x=x^0,y=y^0,z=z^0\) etc. \(\xi=-[al],\ \eta=-[bl],\ \zeta=-[cl]\) etc., respondebit valor ipsius \(T\) hic
\[
T=[ll]-\frac1P
\]
quod vel per expressionem tertiam ponderis \(P\) cum aequatione (16) comparatam, vel per quartam sponte elucet; posteriori, in quo \(x=A,y=B,z=C\) etc., atque \(\xi=\mathfrak A,\ \eta=\mathfrak B,\ \zeta=\mathfrak C\) etc., respondet valor \(T=S\), uti vel e formulis (10) et (15), vel ex his (14) et (16) manifestum est.

\Article{13}
Iam negotium principale consistit in transformatione functionis \(T\) ei simili, quam in Theoria Motus Corporum Coelestium art. 182 atque fusius in Disquisitione de elementis ellipticis Palladis exposuimus. Scilicet statuemus~(17)
\[
\begin{aligned}\relax
[bb,1]&=[bb]-\frac{[ab]^2}{[aa]}\\[4pt]
[bc,1]&=[bc]-\frac{[ab][ac]}{[aa]}\\[4pt]
[bd,1]&=[bd]-\frac{[ab][ad]}{[aa]}\\
\text{etc.}&\\
[cc,2]&=[cc]-\frac{[ac]^2}{[aa]}-\frac{[bc,1]^2}{[bb,1]}\\[4pt]
[cd,2]&=[cd]-\frac{[ac][ad]}{[aa]}-\frac{[bc,1][bd,1]}{[bb,1]}\\
\text{etc.}&\\
[dd,3]&=[dd]-\frac{[ad]^2}{[aa]}-\frac{[bd,1]^2}{[bb,1]}-\frac{[cd,2]^2}{[cc,2]}
\end{aligned}
\]
\SourceLeaf{075}
etc. etc. Dein statuendo\footnote{In praecedentibus sufficere poterant ternae literae pro variis systematibus quantitatum ad tres primas aequationes conditionales referendae: hoc vero loco, ut algorithmi lex clarius eluceat, quartam adiungere visum est; et quum in serie naturali literas \(a,b,c;\ A,B,C;\ \mathfrak A,\mathfrak B,\mathfrak C\) sponte sequantur \(d,D,\mathfrak D\) in serie \(x,y,z\), deficiente alphabeto, apposuimus \(w\), nec non in hac \(\xi,\eta,\zeta\) hanc \(\varphi\).}
\[
\begin{aligned}\relax
[bb,1]y+[bc,1]z+[bd,1]w+\text{etc.}&=\eta'\\
[cc,2]z.+[cd,2]w+\text{etc.}&=\zeta''\\
[dd,3]w+\text{etc.}&=\varphi'''
\end{aligned}
\]
\begin{center}etc., erit\end{center}
\[
T=\frac{\xi\xi}{[aa]}+\frac{\eta'\eta'}{[bb,1]}+\frac{\zeta''\zeta''}{[cc,2]}+\frac{\varphi'''\varphi'''}{[dd,3]}+\text{etc.}
\]
quantitatesque \(\eta',\zeta'',\varphi'''\) etc. a \(\xi,\eta,\zeta,\varphi\) etc. pendent per aequationes sequentes:
\[
\begin{aligned}
\eta'&=\eta-\frac{[ab]}{[aa]}\xi\\[5pt]
\zeta''&=\zeta-\frac{[ac]}{[aa]}\xi-\frac{[bc,1]}{[bb,1]}\eta'\\[5pt]
\varphi'''&=\varphi-\frac{[ad]}{[aa]}\xi-\frac{[bd,1]}{[bb,1]}\eta'-\frac{[cd,2]}{[cc,2]}\zeta''\\
\text{etc.}&
\end{aligned}
\]

Facile iam omnes formulae ad propositum nostrum necessariae hinc desumuntur. Scilicet ad determinationem correlatorum \(A,B,C\) etc. statuemus (18)
\[
\begin{aligned}
\mathfrak B'&=\mathfrak B-\frac{[ab]}{[aa]}\mathfrak A\\[5pt]
\mathfrak C''&=\mathfrak C-\frac{[ac]}{[aa]}\mathfrak A-\frac{[bc,1]}{[bb,1]}\mathfrak B'\\[5pt]
\mathfrak D'''&=\mathfrak D-\frac{[ad]}{[aa]}\mathfrak A-\frac{[bd,1]}{[bb,1]}\mathfrak B'-\frac{[cd,2]}{[cc,2]}\mathfrak C''
\end{aligned}
\]
etc., ac dein \(A,B,C,D\) etc. eruentur per formulas sequentes, et quidem ordine inverso, incipiendo ab ultima,
\[
\left.\begin{aligned}\relax
[aa]A+[ab]B+[ac]C+[ad]D+\text{etc.}&=\mathfrak A\\
[bb,1]B+[bc,1]C+[bd,1]D+\text{etc.}&=\mathfrak B'\\
[cc,2]C+[cd,2]D+\text{etc.}&=\mathfrak C''\\
[dd,3]D+\text{etc.}&=\mathfrak D'''
\end{aligned}\right\}\ (19)
\]
\begin{center}etc.\end{center}
